\documentclass[12pt]{article}
\usepackage[utf8]{inputenc}
\usepackage[T1]{fontenc}
\usepackage{hyperref}
\title{Dongxiang Language Resources for K-12 Curriculum Development}
\author{Compiled by Claude Sonnet 5.5 (claude-sonnet-5-5) for Andr\'as Kornai}
\date{}
\begin{document}
\maketitle

\section{Introduction}
Dongxiang (also called Santa; ISO 639-3 sce) is a Mongolic language of Gansu, China. It has about 200,000 speakers. UNESCO classifies it as vulnerable. There is no long-standing standard script. An older Arabic-script tradition of writing it exists, and a Latin-based orthography has been in trial use since the early 2000s.

\section{Primary Educational Resources}
\begin{itemize}
\item Bilingual teaching experiment at Naleisi Primary School, Dongxiang Autonomous County. A Dongxiang-Chinese bilingual class ran from about 2002 to 2010 using the textbook series \emph{Dongxiangyu xiao gushi} (Dongxiang short stories), whose 8th volume was written by Ma Xinglong and Wang Yilin. The search results mention the textbooks but give no publisher details. See the Lefort article: \url{https://www.jbe-platform.com/content/journals/10.1075/lali.00071.lef}
\item Reports of a 2004 Ford Foundation bilingual pilot: \url{https://pinyin.info/news/wp-json/wp/v2/posts/324}
\end{itemize}

\section{Linguistic and Reference Materials}
\begin{itemize}
\item Santa language (Wikipedia): \url{https://en.wikipedia.org/wiki/Santa_language}
\item Liu Zhaoxiong (1981), \emph{Dongxiangyu jianzhi} (sketch grammar), listed on CrossGram: \url{https://crossgram.clld.org/sources/siniticbetween-Liu1981}
\item Ma Guozhong and Chen Yuanlong (2001), Dongxiang-Chinese dictionary (Gansu Nationalities Publishing House), listed in the Endangered Languages Project bibliography: \url{https://endangeredlanguages.com/lang/1209/bibliography}
\item Lefort, ``Rethinking Dongxiang literacy: the tuhua orou writing system'' (Language and Linguistics): \url{https://benjamins.com/catalog/lali.00071.lef} and \url{https://www.altaist.org/home/confessions/lefort-julie/rethinking-dongxiang-literacy-the-tuhua-orou-writing-system-and-its-sociolinguistic-practices/}
\item Omniglot chart for Santa: \url{https://www.omniglot.com/charts/santa.pdf}
\item Kim, ``Santa,'' in \emph{The Mongolic Languages} (Routledge, 2003), cited in the search results but with no online text found.
\end{itemize}

\section{Religious and Community Resources}
No Dongxiang scripture or religious texts were found. Many Dongxiang are Muslim, and the community writes the language informally in Arabic script (tuhua orou). Further resources: \url{https://rosettapanglossia.longnow.org/index.php?title=Sce}

\section{Digital and Audio Resources}
\begin{itemize}
\item Dongxiang Short Stories (with gloss and annotations), Vol. 1, on HAL-INALCO: \url{https://hal-inalco.archives-ouvertes.fr/hal-04015834v2}
\item Freelang Dongxiang-English word list (606 Dongxiang-English and 481 English-Dongxiang entries): \url{https://www.freelang.net/dictionary/dongxiang.php}
\item Chinese-to-Dongxiang machine translation demo (fine-tuned NLLB model): \url{https://dxlanguage-chinesedongxiangtranslation.hf.space}
\item Numerals data (MPI-EVA): \url{https://lingweb.eva.mpg.de/channumerals/Dongxiang.htm}
\end{itemize}

\section{Curriculum Development Context}
Before the early 2000s, Dongxiang had no standard written form, and only a small amount of schooling was done in it. The experimental bilingual classes (2002--2010) produced primary-level readers. Their current status is unclear. A broader curriculum would first need a settled orthography, a choice between Latin and Arabic script, a modern dictionary and graded readers.

\end{document}